\chapter{Register Refactoring}
\label{ch:registers}


\section{Register Types}

Fiat Cryptography was originally designed around scalar x86 registers, and therefore had no notion of registers larger than 64 bits. AVX includes 128-bit xmm and 256-bit ymm registers, so we needed to refactor the entire register system to accomodate these. 
The simplest change began with the register type. Scalar and vector registers are fundamentally different in terms of how they interact with the rest of the codebase, in a few important ways:
\begin{itemize}
    \item Only scalar registers can form memory-address operands
    \item Only scalar registers can participate in flag-setting arithmetic and carry chains
    \item Only scalar registers can appear in calling conventions
    \item Scalar registers' alias hierarchy is four levels deep (al/ax/eax/rax)
    \item Vector registers carry opaque packed-lane data
    \item Vector registers alias only two levels deep (xmm/ymm)
\end{itemize} we made the two register kinds distinguishable by type, with a unifying type REG = SReg SREG | VReg VREG.

\codeshot{REG.png}

Downstream functions which need to distinguish between the two can now do so directly by type. For example, a memory location needs to accept only scalar registers for the base address:

\codeshot{MEM.png}


\section{Register get/set}

Many existing functions for interacting with register values in Fiat-Crypto were specialized to handling 64-bit values. For example, the SetReg function uses two branches for 64-bit writes and smaller register writes because an 8 or 32-bit write neesd to combine the new value iwth the existing value, whereas a 64-bit write can overwrite hte old value, meaning the old value can be completely forgotten (or be unspecified to begin with), simplifying the computation graph in many cases.

\codeshot{SetRegOld.png}

However, when writing to a 128-bit wide xmm register, even a 64-bit write needs to set a slice of the old value, so this function should check more generally if the write matches the largest size of this register's aliasing hierarchy:

\codeshot{SetRegNew.png}

This kind of refactor requires new lemma structure as well. In this case, we introduced two new lemmas:

\codeshot{R_SetReg_full.png}
\vspace{-5ex}
\codeshot{R_SetReg_partial.png}
